\cotefile{78}
\scanpage{75}
\noindent\begin{tabular}{||p{.9\linewidth}}\unsure{Réalisation|Réalisations} géométriques des polyèdres \deleted{\unsure{infinis|affines}} combinatoires
en coordonnées affines universelles.\end{tabular}

\uline{1}. Programme de travail.

\deleted{\unsure{1.}} \uline{1}. Soit $\Pi$ polyèdre combinatoire donné, on se
propose d'étudier ses réalisations géométriques \inserted{régulières}
dans des \deleted{\unsure{espaces}} \inserted{plans} affines variables $E$. Notons
que si deux telles réalisations sont $\Pi$-isom, il y
a un $\Pi$-isom \uline{unique} entre elles, par
définition. Donc dans tte classe d'isomorphie
(\unsure{notée} $\alpha$, pour fixer les idées), il y a un
représentant \inserted{$(E_\alpha,\varphi_{\alpha s},\varphi_{\alpha\unsure{d|a}})$}, défini à isom unique près.
En particulier, le plan affine \inserted{$E_\alpha$} de cette réalisation
géom est défini à \uline{isom. unique} \uline{près}.
On indexera $\alpha$ à l'aide d'un invariant
numérique (noté également $\alpha$ -- \unsure{essentiellt|essentiellement|en icelles|ensemble}
ce sera le \inserted{``\unsure{inv. numérique}''} $\unsure{\zeta+\zeta^{-1}|\zeta+\zeta'}$ déjà envisagé précédemment)
et pour bien faire il faudra pouvoir
expliciter $E_\alpha,\varphi_{\alpha s},\varphi_{\alpha d}$ \inserted{(à isom unique près)} à l'aide de $\Pi$ et
de cet inv. numérique.

On va pour commencer se donner
un repère $\overline{\bigl|\,r\in\operatorname{Rep}(\Pi)\,\bigr|}$
pour épingler $\Pi$, de sorte que $\Pi$ devient
% archive sheet number: 74
